\clearpage
\VerdierSemanticAnchor{front-editors-note}
\VerdierPhysicalPageAnchor{0008}
\chapter*{编者按}
\addcontentsline{toc}{chapter}{编者按}

\VerdierSemanticAnchor{front-editors-note-publication-history}
1989 年 10 月 19 日，Université Paris VII 举行纪念 J.-L. Verdier 的仪式；
M. Artin、P. Cartier 和 L. Illusie 正是在那里萌生了出版 J.-L. Verdier
博士论文的想法。两年后，我加入他们，并在 J. Le Potier 与 M. Broué 的支持
下，从 Paris VII（URA 212）和 ÉNS（DMI）获得经费，将本文录入 \TeX{}。
录入完成后，我们把结果与 Verdier 的打字原稿作了仔细核对。由于 Verdier
从未作过最后修订，仍须进行一次数学内容上的审读。

\VerdierSemanticAnchor{front-editors-note-chapter-i-modifications}
这项细致的审读工作由 P. Cartier（第一章）、B. Keller（第三章）以及
M. Zisman 和我（全文）极其认真地完成。除大量不值得逐一列举的排印错误外，
主要改动如下。第一章命题 (3.1.5) 原来断言一个等价关系。我们把等价改为蕴含，
因为其逆命题（后文从未使用）是错误的。命题 (3.4.1) 有一个符号错误，并由此
造成第三章 ((1.3.4)、(b) 及 (3.2.7.3)) 的若干符号错误。这又要求调整“对偶
三角范畴”（第二章，1.1.7）的定义，以使下述断言仍然成立：\emph{对偶范畴的
导出范畴，就是原导出范畴的对偶三角范畴}。

\VerdierSemanticAnchor{front-editors-note-chapter-ii-modifications}
第二章中，§1.2 各命题的次序有所调整。Verdier 写完正文后认识到，两个特异
三角的直和仍为特异三角这一断言，可由三角范畴的其他公理推出。因此他把它从
公理表中删去，却没有完成由此产生的重新编排。此外，后文未使用且验证失败的
(2.3.6) 和 (2.3.7) 两段也被删除。定常谱对象 (4.4.2) 的定义作了轻微修改，
因为原定义过弱，不足以证明随后的定理 (4.4.3)。我们核实了这一修改不会影响
后续内容。

\clearpage
\VerdierPhysicalPageAnchor{0009}
\VerdierSemanticAnchor{front-editors-note-apparatus}
记号索引和术语索引都得到了补充。与此相反，参考文献保持不变（但 SGA 各卷
改列其在 ``Lecture Notes'' 中的最终书目信息）。尤其是，参考文献编号 13 的
缺失并未补上。

\VerdierSemanticAnchor{front-editors-note-acknowledgements}
谨向 M. Broué、A. Bruguières、P. Cartier、L. Gruson、L. Illusie、
B. Keller、J. Le Potier、M. Zisman，以及 Université Paris VII、École
Normale Supérieure 和 Institut des Hautes Études Scientifiques 致以
诚挚谢意；他们以不同方式帮助本项目得以完成。我还感谢 C. Defosse 和
C. Gourgues 耐心地将本书录入 \textsc{AMS}-\TeX{}。

\VerdierSemanticAnchor{front-editors-note-memorial}
我本也希望能向 L. Doustaing 表达同样的谢意；遗憾的是，他已于 1994 年
12 月英年早逝。

\VerdierSemanticAnchor{front-editors-note-signature}
\begin{flushright}
巴黎，1996 年 6 月 14 日\\
Georges Maltsiniotis
\end{flushright}
